% Pandas & Matplotlib Cheatsheet
\scriptsize
\newgeometry{margin=.5cm}
\lstset{style=python, basicstyle=\ttfamily\scriptsize}
\pagestyle{empty}
\raggedcolumns
\footnotesize
\newcommand{\cheatSheetColor}{ForestGreen}
\newtcolorbox{cheatSheetBox}[1][]{
  enhanced jigsaw,
  colback=\cheatSheetColor!10,
  colframe=\cheatSheetColor,
  colbacktitle=\cheatSheetColor!30,
  fonttitle=\bfseries,
  parbox=false,
  attach boxed title to top text left={yshift=-5mm,yshifttext=-3.5mm},
  coltitle=black,
  #1}

\begin{landscape}
  \footnotesize
  \begin{multicols*}{3}
    \textcolor{\cheatSheetColor}{\begin{center}
        \section*{Data Science Cheatsheet}
      \end{center}\hrule
      Das folgende Cheatsheet enthält die wichtigsten Befehle für den Umgang mit den Libraries \texttt{pandas}, \texttt{matplotlib} und \texttt{statsmodels}.
    }

    % Pandas Basics
    \begin{cheatSheetBox}[title=Libraries Importieren]
      \begin{lstlisting}
import pandas as pd
import matplotlib.pyplot as plt
import statsmodels.formula.api as smf
\end{lstlisting}
    \end{cheatSheetBox}

    \begin{cheatSheetBox}[title=CSV-Datei laden]
      \begin{lstlisting}
df = pd.read_csv('datei.csv')
\end{lstlisting}
    \end{cheatSheetBox}

    \begin{cheatSheetBox}[title=Erste/letzte Zeilen anzeigen]
      \begin{lstlisting}
df.head()   # erste 5 Zeilen
df.tail()   # letzte 5 Zeilen
\end{lstlisting}
    \end{cheatSheetBox}

    \begin{cheatSheetBox}[title=Typ einer Variable anzeigen]
      \begin{lstlisting}
type(df) # Tabelle (DataFrame)
type(df['col']) # Spalte (Series)
type(1) # Zahl (int)
type('Text') # Text (str)
df.dtypes # Daten-Typen aller Spalten
\end{lstlisting}
    \end{cheatSheetBox}

    \begin{cheatSheetBox}[title=Spalten auswählen]
      \begin{lstlisting}
df['Spalte']   # einzelne Spalte
df[['A', 'B']] # mehrere Spalten
\end{lstlisting}
    \end{cheatSheetBox}

    \begin{cheatSheetBox}[title=Aggregatsfunktionen]
      \begin{lstlisting}
df['Spalte'].mean() # Mittelwert
df['Spalte'].sum()  # Summe
df['Spalte'].min()  # Minimum
df['Spalte'].max()  # Maximum
df['Spalte'].count() # Anzahl Einträge
df['Spalte'].std()  # Standardabweichung
df['Spalte'].unique() # Einzigartige Werte
df['Spalte'].value_counts() # Häufigkeit der Werte
\end{lstlisting}

    \end{cheatSheetBox}

    \begin{cheatSheetBox}[title=Zeilen filtern]
      \begin{lstlisting}
df[df['Alter'] > 18] # Filter 
\end{lstlisting}
    \end{cheatSheetBox}

    \begin{cheatSheetBox}[title=Neue Spalte erstellen]
      \begin{lstlisting}
df['BMI'] = df['Gewicht'] / (df['Grösse']/100)**2
\end{lstlisting}
    \end{cheatSheetBox}

    \begin{cheatSheetBox}[title=Grafiken erstellen]
      \begin{lstlisting}
df["column"].value_counts().plot.pie() # Kuchendiagramm
df["column"].value_counts().plot.bar() # Balkendiagramm
df.plot.line(x='x_column', y='y_column') # Liniendiagramm
df.plot.scatter(x='x_column', y='y_column') # Streudiagramm

# scatter plots nach Kategorien einfärben
df['category_column'] = df['category_column'].astype('category') 
df.plot.scatter(x='x_column', y='y_column', c='category_column', colormap='viridis')

# Histogramm mit 10 Grössenklassen
df['column'].plot.hist(bins=10)
\end{lstlisting}
    \end{cheatSheetBox}

    \begin{cheatSheetBox}[title=Gruppieren]
      \begin{lstlisting}
df.groupby('Spalte').mean()
\end{lstlisting}
    \end{cheatSheetBox}

    %\section*{Lineare Regression}

    \begin{cheatSheetBox}[title=Lineare Funktion]
      Formel einer linearen Funktion:
      \[y = m \cdot x + q\]
      Wobei m die Steigung und q den y-Achsenabschnitt darstellt.
      \[m=\frac{y_2 - y_1}{x_2 - x_1}\]
      \[q = y_1 - m \cdot x_1\]
    \end{cheatSheetBox}
    \begin{cheatSheetBox}[title=Lineare Regression]
      \begin{itemize}
        \item \textbf{Unabhängige Variable}: \lstinline|df['x']|
        \item \textbf{Abhängige Variable}: \lstinline|df['y']|
      \end{itemize}
      \begin{lstlisting}
model = smf.ols(formula='y ~ x', data=df).fit()
print(model.params) # Koeffizienten

# Vorhersage
df['y_pred'] = model.predict(df['x'])
\end{lstlisting}

    \end{cheatSheetBox}
  \end{multicols*}
\end{landscape}
